% ECONOMICS_PROBLEM: {"id": "C73-77", "title": "Pure positional equilibria in multiplayer reachability games", "jel_code": "C73", "source_id": "OP-0161"}
% !TeX program = xelatex
\documentclass[11pt,a4paper]{article}
\usepackage[margin=23mm]{geometry}
\usepackage{fontspec}
\setmainfont{Latin Modern Roman}
\usepackage{amsmath,amssymb,mathtools}
\usepackage{xcolor,enumitem,booktabs,longtable,array,xurl,hyperref}
\definecolor{muted}{HTML}{53636C}
\hypersetup{colorlinks=true,urlcolor=blue,linkcolor=blue}
\setlength{\parindent}{0pt}
\setlength{\parskip}{5pt}
\newcommand{\R}{\mathbb R}
\newcommand{\E}{\mathbb E}
\newcommand{\N}{\mathbb N}
\newcommand{\ind}{\mathbf 1}
\newcommand{\Var}{\operatorname{Var}}
\newcommand{\Cov}{\operatorname{Cov}}
\newcommand{\supp}{\operatorname{supp}}
\DeclareMathOperator*{\argmin}{arg\,min}
\DeclareMathOperator*{\argmax}{arg\,max}
\newcommand{\entrymeta}[3]{{\sffamily\small\color{muted}\textbf{\texttt{#1}}\quad|\quad #2\par #3\par}}
\newcommand{\Problemref}[1]{\texttt{#1}}
\newcommand{\JELref}[2]{\texttt{#2} (JEL \texttt{#1})}
\begin{document}
\section*{Pure positional equilibria in multiplayer reachability games}
\textbf{Problem ID:} \texttt{C73-77}\quad \textbf{JEL:} \texttt{C73}\par
% BEGIN ECONOMICS STATEMENT
\label{problem:OP-0161}
% GAME_THEORY_PROBLEM: {"local_id": "GT-031", "original_number": 31, "kind": "P", "entry_kind": "statement_only", "published": false, "review_required": true, "category": "game_theory"}
\label{game:GT-031}
{\sffamily\small\color{muted}
{\sffamily\small\color{muted}\textbf{GT-031} | Qualitative and graph games\par P | Source-reported pure positional existence question\par}
\par}
\subsection*{Definitions and mathematical statement}
Let $(V,E)$ be a finite directed graph with at least one outgoing edge at every vertex, and let $V=\bigsqcup_{i=1}^{n}V_i$ specify its controlling players. There are no stochastic vertices. Each player has a target $T_i\subseteq V$. A behavioral strategy chooses a distribution over successors at every finite vertex history ending in $V_i$. A pure positional strategy is a function $f_i:V_i\to V$ with $(v,f_i(v))\in E$. For initial vertex $v_0$, the payoff is
\[
 u_i^{v_0}(\sigma)=\Pr_{v_0,\sigma}(\exists t\geq0:\ v_t\in T_i).
\]
The question is
\[
 \forall n\geq2\ \forall(V,E,(V_i),(T_i))\ \forall v_0\in V
 \quad\exists(f_i)_{i=1}^{n}\quad\forall i\ \forall\tau_i,
 \qquad u_i^{v_0}(\tau_i,f_{-i})\leq u_i^{v_0}(f),
\]
where deviations may be arbitrary randomized history-dependent strategies.
\subsection*{Short English statement}
Does every finite deterministic turn-based multiplayer reachability game have a Nash equilibrium using only the current vertex and no randomization?
\subsection*{Variants and scope boundaries}
An everywhere Nash equilibrium uses one profile satisfying the Nash inequalities from every vertex, with each objective started afresh. A positional SPE must satisfy the continuation inequalities after every feasible history and retain earlier target visits. These are separate stronger targets. Randomized memoryless strategies are not positional in the source's terminology. Absorbing targets form a special positive subcase, not the unrestricted model.
\subsection*{Source relationship and status}
[S1], Section 5, expressly leaves reachability-only pure positional NE and stronger refinements open. Its positive results for randomized stationary equilibria do not settle this question. Theorem 10 proves a positional SPE under additional absorption assumptions.
\subsection*{References}
\begingroup\footnotesize\raggedright
\textbf{[S1]} Mona Alluwaym, James C. A. Main and Sven Schewe, \emph{Simple Nash Equilibria for Qualitative Multiplayer Games}, MFCS 2026, LIPIcs 386, Article 85, pp. 85:1--85:17, doi:10.4230/LIPIcs.MFCS.2026.85. Locators: Section 2, pp. 85:4--85:5 (strategy and refinement definitions); Theorem 10 and Section 5, p. 85:14 (special absorption result and open positional cases). \url{https://drops.dagstuhl.de/storage/00lipics/lipics-vol386-mfcs2026/LIPIcs.MFCS.2026.85/LIPIcs.MFCS.2026.85.pdf}
\par\endgroup

\subsection*{Common conventions: Source mathematical conventions}

\textbf{Finite games.} For a finite set $A$, $\Delta(A)$ is its probability
simplex. Players choose independent mixed strategies in Nash equilibrium;
correlation is allowed only when explicitly part of the solution concept.
In a game with payoff functions $u_i$ and strategy profile $x$, an additive
$\varepsilon$ Nash equilibrium satisfies
\[
 u_i(x)\geq u_i(x_i',x_{-i})-\varepsilon
 \quad\text{for every player $i$ and every admissible deviation $x_i'$}.
\]
For numerical approximation, payoffs are normalized as locally specified.
An $\varepsilon$ payoff residual is distinct from distance at most
$\varepsilon$ to the set of exact equilibria. Point convergence, convergence
of empirical play, and convergence in distance to an equilibrium set are
also distinct.

\textbf{Computational representations.} Unless stated otherwise, a complexity
question takes rational data with binary encoding; polynomial time is measured
in total bit length. A fixed-accuracy polynomial algorithm is not necessarily
polynomial in $1/\varepsilon$, and neither is necessarily polynomial in
$\log(1/\varepsilon)$. Succinct games and value, demand, or sampling oracles
must declare their interfaces. Exact real solutions require an explicit
representation; an unspecified list of arbitrary real numbers is not a
finite computational output. A claimed algorithmic barrier is meaningful
only for its specified model and reductions.

\textbf{Dynamic games.} Histories, information, observation, and allowable
randomization are part of the model. A stationary strategy depends only on
the current state; a positional strategy in a deterministic graph game
chooses one action at each controlled vertex. Nash equilibrium at an initial
state and equilibrium after every history are different requirements.
An expectation of a pathwise $\liminf$ payoff, a $\liminf$ of expected averages,
a limiting expected average, and a uniform finite-horizon equilibrium are
different criteria. None is silently substituted for another.

\textbf{Allocation and mechanisms.} A complete allocation assigns every item
exactly once unless otherwise stated. Zero-valued goods, negative values,
capacity constraints, feasibility, lotteries, and copies of goods affect
fairness definitions and are specified locally. Dominant-strategy incentive
compatibility, Bayesian incentive compatibility, universal truthfulness,
and truthfulness in expectation impose different inequalities. Whenever
an objective ratio has zero denominator, its convention is stated or those
instances are excluded.

\textbf{Infinite-dimensional models.} Expectations, integrals, stochastic
equations, and optimization objectives require their local regularity,
integrability, filtrations, and admissible domains. A backward--forward
mean-field system is a boundary-value problem; it is not an initial-value
semiflow without an additional construction. Long-time, large-population,
and vanishing-noise limits require an order of limits and a convergence
topology. If a minimum need not be attained, the objective is its infimum
or supremum and the approximation requirement is stated separately.
% END ECONOMICS STATEMENT
\end{document}
